%% LyX 2.2.3 created this file.  For more info, see http://www.lyx.org/.
%% Do not edit unless you really know what you are doing.
\documentclass[11pt,ngerman,intoc]{article}
\usepackage[T1]{fontenc}
\usepackage[latin9]{inputenc}
\usepackage{geometry}
\geometry{verbose,tmargin=3cm,bmargin=2cm,lmargin=2cm,rmargin=2.3cm}
\setcounter{secnumdepth}{5}
\setcounter{tocdepth}{5}
\usepackage{float}
\usepackage{amsmath}
\usepackage{setspace}
\onehalfspacing

\makeatletter
%%%%%%%%%%%%%%%%%%%%%%%%%%%%%% Textclass specific LaTeX commands.
\numberwithin{equation}{section}
\numberwithin{table}{section}
\numberwithin{figure}{section}

%%%%%%%%%%%%%%%%%%%%%%%%%%%%%% User specified LaTeX commands.
\RequirePackage[ngerman=ngerman-x-latest]{hyphsubst}
\usepackage{caption}
\usepackage{lmodern}
\usepackage[bottom]{footmisc}
\usepackage{setspace}
\usepackage[natbibapa]{apacite} 
\onehalfspacing
\hyphenation{Dreh-mo-men-ten-gleich-ge-wicht} 
\newtheorem{defi}{Definition}[section]
\usepackage{url}
\usepackage{ucs}
\PrerenderUnicode{������}
\makeatletter
\newcommand{\tabitem}{~~\llap{\textbullet}~~}
\usepackage{pdfpages}
\usepackage{acronym}
\usepackage{graphicx}
\usepackage[colorlinks=false,
pdfpagelabels,
pdfstartview = FitH,
bookmarksopen = true,
bookmarksnumbered = true,
linkcolor = black,
plainpages = false,
hypertexnames = false,
citecolor = black] {hyperref}
\setlength{\parindent}{0pt}
\PassOptionsToPackage{hyphens}{url}
\usepackage{hyperref}
\expandafter\def\expandafter\UrlBreaks\expandafter{\UrlBreaks% save the current one
  \do\a\do\b\do\c\do\d\do\e\do\f\do\g\do\h\do\i\do\j%
  \do\k\do\l\do\m\do\n\do\o\do\p\do\q\do\r\do\s\do\t%
  \do\u\do\v\do\w\do\x\do\y\do\z\do\A\do\B\do\C\do\D%
  \do\E\do\F\do\G\do\H\do\I\do\J\do\K\do\L\do\M\do\N%
  \do\O\do\P\do\Q\do\R\do\S\do\T\do\U\do\V\do\W\do\X%
  \do\Y\do\Z\do\*\do\-\do\~\do\'\do\"\do\-}%

\makeatother

\usepackage{babel}
\begin{document}
\includepdf{Deckblatt.pdf}
\setcounter{page}{2}

\section{Klassendiagramme (f�r den Problembereich)}
Unser Hauptproblem war, die Komplexit�t des Diagramms einzuschr�nken. Dies zeigte sich darin, dass zwar Aggregation und Komposition herausgelassen werden, allerdings Typen und Visibilities einbezogen wurden. Au�erdem wurden 5 der F�hrungspositionen im Soll-Diagramm dargestellt, was ebenfalls nicht n�tig gewesen w�re.\\
Gut geklappt hat die Darstellung der Daten der Helfer. Trotz anf�nglicher Unklarheiten, wie die Daten strukturiert werden sollten, konnten wir eine gute L�sung f�r dieses Problem finden, die es erlaubte verschiedene Stammdaten und Qualifikationen eines Helfers besonders kompakt darzustellen. Dabei waren vor allem die gegeben Excel-Tabellen der Malteser hilfreich.
\section{Aktivit�tsdiagramme (f�r Gesch�ftsprozesse)}
An sich halten wir Aktivit�tsdiagramme f�r geeignet f�r Gesch�ftsprozesse, jedoch h�tten wir die UML hier gern erweitert. Uns fehlte eine Aktivit�t, welche den Postausgang eines Mitgliedsantrags symbolisiert, sowie eine M�glichkeit, eindeutig darzustellen, dass auf eine Antwort reagiert wird.\\
Wir h�tten gern den Unterschied zwischen Aktionen und Aktivit�ten gewusst, da wir so die Verschachtelung der Gesch�ftsprozesse besser h�tten darstellen k�nnen.\\
Unser Hauptproblem bei der Gesch�ftsprozessmodellierung war jedoch, dass wir bei den Interviewterminen scheinbar leider nicht alles richtig verstanden und deswegen bestimmte Details falsch dargestellt haben, etwa den Prozess der Mitgliedsaufnahme.
\section{Anwendungsfalldiagramme}
\section{Prozess der objektorientierten Analyse}
Bei der Analyse ist in unserer Gruppe positiv anzumerken, dass es uns gelungen ist die einzelnen Akteure und Objekte des Problembereichs klar auszuarbeiten und anschlie�end deren Relationen und Gesch�ftsprozesse darzustellen.
Ausbauf�hig ist bei uns jedoch die Konsistenz der Benamung verschiedener Objekte zwischen verschiedenen Diagrammen.
Weiterhin l�sst sich anmerken, dass wir den Problembereich mehr als Datenmodell dargestellt haben, als das komplette System mit allen Akteuren zu modellieren.\\
\section{Modellierungswerkzeuge}
\section{Interviewf�hrung}
\section{Teamwork und Kommunikation}
Obwohl jede Aufgabe im Team und nicht als Einzelarbeit bearbeitet wurde, gab es nie hitzig emotionale Diskussionen, sodass Konflikte immer schnell und sachlich gel�st werden konnten, was eine gute Zusammenarbeit garantiert hat.\\
Die Kommunikationsmittel waren ein HackMD, ein Slack-Channel und eine Telegram-Gruppe. �ber letztere lief die Organisation der Treffen, die anderen beiden waren zu Diskussion und Dokumentation von Ergebnissen gedacht.\\ Allerdings war gerade in diesen die Response Time deutlich geringer und die Antworten auch unorganisierter.\\
Letztendlich h�tte man den Tutor st�rker einbinden sollen und die Treffen konsequenter organisieren m�ssen.\\

\end{document}
